\chapter{Virtual Ecology}

\chapterepigraph{%
  A simulation that does not surprise you\\
  is not a simulation.\\
  It is a description.
}{Flyxion, \textit{Semantic Infrastructure}}

\sidenote{Connects to\\constraint ecology\\(Ch.\,26) and\\accessibility\\fields (App.\,J)}

Chapter~41 developed procedural generation as constraint satisfaction.
This chapter examines the emergent communities — virtual ecologies — that
arise in richly constrained simulated environments, and what they reveal
about real ecologies.

\section{Agents and Resources}

A virtual ecology begins with two primitive elements: agents (entities
that act) and resources (quantities that agents consume and produce).
The constraint structure governing their interaction determines what
ecologies emerge.

\begin{definition}{Virtual Ecology}{virtual-ecology}
  A \emph{virtual ecology} is a dynamical system $(A, R, E, \kappa_E)$
  where:
  \begin{itemize}
    \item $A = \{a_1, \ldots, a_n\}$ is a set of agents with strategies
          $\pi_i : \mathcal{O}_i \to \mathcal{A}_i$,
    \item $R = \{r_1, \ldots, r_m\}$ is a set of resources,
    \item $E : A \times R \to \mathbb{R}$ is the exchange matrix
          (how much of each resource each agent type consumes and
          produces),
    \item $\kappa_E$ is the ecological constraint functional
          (resource levels must remain non-negative; population
          levels must be non-negative).
  \end{itemize}
\end{definition}

\section{Adaptation}

In biological ecologies, adaptation occurs through evolution: agents
with higher-fitness strategies reproduce more. In virtual ecologies,
adaptation can be implemented directly: agent strategies update in
response to outcomes.

The accessibility framework gives a precise interpretation: adaptation
is constraint descent in the strategy space. Agents move toward
strategies that minimize their constraint violation (maximize survival
and reproduction) given the current state of the ecology.

\section{Civilizations}

When virtual agents develop complex strategies — division of labor,
trade, coordination, conflict — virtual civilizations emerge. The
transition from ecology to civilization is a phase transition in the
accessibility landscape: new basins appear (stable social structures)
that were not present in the pre-civilizational ecology.

\begin{namedthm}{The Civilization Emergence Theorem (Virtual)}
  In a virtual ecology with sufficient agent capability (complex
  strategies) and resource diversity (multiple resource types with
  complementary advantages), civilization-like structures emerge
  as accessibility attractors: stable configurations where specialization
  and exchange reduce total constraint violation below what any single
  agent strategy can achieve.
\end{namedthm}

\begin{exercise}
  Describe the minimal virtual ecology that produces civilization-like
  trade. What resource types are needed? What agent capabilities?
  What constraint structure produces the emergence of specialization?
\end{exercise}

\begin{exercise}
  Many virtual civilizations in games (Civilization, Crusader Kings,
  Dwarf Fortress) exhibit ``emergent stories'' — narratives that arise
  from the interaction of rules rather than being scripted. Using
  the accessibility framework, explain what makes a game rich enough
  to produce emergent stories. What structural property distinguishes
  story-rich from story-poor games?
\end{exercise}
